% 全部作者、卷期及文章号见docs/biography-sources.md；作者身份由申请人确认。
% 本节为完整成果目录，不采用正文文献表的“et al.”截断规则。
\noindent\textbf{代表性研究成果}

\begingroup
\fontsize{10.5pt}{16bp}\selectfont\frenchspacing
\begin{enumerate}[label=(\arabic*),leftmargin=2.2em,labelsep=.5em,topsep=5bp,itemsep=6bp,parsep=0pt]
  \item Xu P, \textbf{Yan F}, Zhou Y, Liu Y, Wu C.
  \href{https://doi.org/10.1016/j.eswa.2026.132953}{RSHR+: Progressive Question-Conditioned Visual Calibration and Structured State-Space Reasoning for Remote Sensing Visual Question Answering}[J].
  \textit{Expert Systems with Applications}, 2026, 328: 132953.\quad\textbf{唯一通讯作者。}

  \item Liu Y, \textbf{Yan F}, Xu P, Liu C, Wang L.
  \href{https://doi.org/10.1016/j.neucom.2026.133361}{WaveKAN: Adaptive wavelet decomposition and Kolmogorov-Arnold network for long-term time series forecasting}[J].
  \textit{Neurocomputing}, 2026, 682: 133361.\quad\textbf{共同第一作者。}

  \item Xu P, \textbf{Yan F}.
  \href{https://doi.org/10.1016/j.cviu.2026.104753}{RSSR: Efficient knowledge transfer and deep spatial modeling for remote sensing visual question answering}[J].
  \textit{Computer Vision and Image Understanding}, 2026, 268: 104753.\quad\textbf{唯一通讯作者。}

  \item Sun Q, \textbf{Yan F}, Sun W, Zhou Y.
  \href{https://doi.org/10.1016/j.energy.2025.139283}{DWT-Former: Fusing wavelet-based multi-scale features and transformer-based temporal representations for photovoltaic power forecasting}[J].
  \textit{Energy}, 2025, 341: 139283.\quad\textbf{唯一通讯作者。}

  \item Wang Z, \textbf{Yan F}, Wang L, Yin Y, Lin J.
  \href{https://doi.org/10.1016/j.neunet.2025.107782}{S-YOLO: An enhanced small object detection method based on adaptive gating strategy and dynamic multi-scale focus module}[J].
  \textit{Neural Networks}, 2025, 191: 107782.\quad\textbf{共同第一作者。}

  \item \textbf{Yan F}, Li Z, Silamu W, Li Y.
  \href{https://doi.org/10.1007/s10994-023-06426-6}{Knowledge-aware image understanding with multi-level visual representation enhancement for visual question answering}[J].
  \textit{Machine Learning}, 2024, 113(6): 3789--3805.\quad\textbf{第一作者、通讯作者。}
\end{enumerate}
\endgroup

\noindent\textbf{本课题承担的任务}

作为课题负责人，负责总体研究设计、技术方案制定和进度管理。重点组织矿区数据来源核验与误差诊断，设计跨时相一致性约束的区域对照实验，分析约束对误报和细小变化漏检的影响；负责概率校准与错误排序方法设计，在固定复核面积下评价预测可靠性。统筹跨矿区验证、矿区级统计和失败案例分析，指导研究生完成算法实现、实验记录与结果复核，并负责论文撰写、学术交流和结题报告编制。
